\documentclass[11pt]{article}

\usepackage[T1]{fontenc}
\usepackage[utf8]{inputenc}
\usepackage{lmodern}
\usepackage[margin=1.06in]{geometry}
\usepackage{amsmath,amssymb,mathtools,amsthm}
\usepackage{aliascnt}
\usepackage{mathrsfs}
\usepackage{microtype}
\usepackage{enumitem}
\usepackage{xcolor}
\usepackage{url}
\usepackage{hyperref}

\hypersetup{
  colorlinks=true,
  linkcolor=blue!45!black,
  citecolor=green!35!black,
  urlcolor=blue!55!black,
  pdftitle={A Linear Kontsevich--Zagier Theorem for Compact Exponential Periods on the Affine Line},
  pdfauthor={Christopher D. Long; Antoine-Auguste Le Blanc},
  pdfsubject={Exponential periods, E-functions, formal periods, polynomial exponential integrals},
  pdfkeywords={exponential period, Kontsevich--Zagier period conjecture, E-function, affine line, twisted de Rham cohomology, Cauchy transform}
}

\numberwithin{equation}{section}
\setlength{\emergencystretch}{3em}

\newtheorem{theorem}{Theorem}[section]

\newaliascnt{proposition}{theorem}
\newtheorem{proposition}[proposition]{Proposition}
\aliascntresetthe{proposition}

\newaliascnt{lemma}{theorem}
\newtheorem{lemma}[lemma]{Lemma}
\aliascntresetthe{lemma}

\newaliascnt{corollary}{theorem}
\newtheorem{corollary}[corollary]{Corollary}
\aliascntresetthe{corollary}

\theoremstyle{definition}
\newaliascnt{definition}{theorem}
\newtheorem{definition}[definition]{Definition}
\aliascntresetthe{definition}

\theoremstyle{remark}
\newaliascnt{remark}{theorem}
\newtheorem{remark}[remark]{Remark}
\aliascntresetthe{remark}

\usepackage[capitalize,noabbrev]{cleveref}

\newcommand{\C}{\mathbb C}
\newcommand{\Qbar}{\overline{\mathbb Q}}
\newcommand{\Aone}{\mathbb A^1}
\newcommand{\Span}{\operatorname{span}}
\newcommand{\supp}{\operatorname{supp}}
\newcommand{\per}{\operatorname{per}}
\newcommand{\Pcal}{\mathcal P}
\newcommand{\Ecal}{\mathcal E}
\newcommand{\Rcal}{\mathcal R}
\newcommand{\Mcal}{\mathcal M}
\newcommand{\Ocal}{\mathcal O}

\title{A Linear Kontsevich--Zagier Theorem for\\
Compact Exponential Periods on the Affine Line}
\author{Christopher D. Long \and Antoine-Auguste Le Blanc\thanks{Le Blanc cryptographic author identifier: \texttt{b6048125\allowbreak{}30b3535b\allowbreak{}23d6dbb0\allowbreak{}f5abe32d\allowbreak{}1ed1a9de\allowbreak{}35fb4f7d\allowbreak{}04d62cff\allowbreak{}08ea19c8}}}
\date{August 2026}

\begin{document}
\maketitle

\begin{center}
\small
\texttt{galizur@gmail.com}
\end{center}

\begin{abstract}
Let $q,r\in\Qbar[x]$, let
$\Delta=\sum_\alpha d_\alpha[\alpha]$ be a finite degree-zero zero-cycle on
$\Aone(\Qbar)$ with coefficients in $\Qbar$, and let $\Gamma$ be any finite
piecewise $C^1$ one-chain in $\C$ with $\partial\Gamma=\Delta$.  We consider
\[
  I_{q,r,\Delta}(z)=\int_\Gamma r(x)e^{zq(x)}\,dx
\]
and the compact exponential period $I_{q,r,\Delta}(1)$.

We define the formal exponential-period space on smooth affine curves using
bilinearity, functoriality, and twisted Stokes, and consider the subspace
generated by the affine-line periods above and by the zero-dimensional exponential periods
$e^\beta$, $\beta\in\Qbar$.  We prove that the period map is injective on
this subspace.  Equivalently, every $\Qbar$-linear numerical relation among
these compact exponential periods is generated by the standard
Kontsevich--Zagier operations.

The proof first reduces every polynomial differential uniquely modulo
$R\mapsto R'+q'R$.  Direct sums of the moment systems from the companion
paper \emph{Algebraic Exponential Integrals} are then used to lift a
numerical relation by Beukers' theorem.  A direct-sum residue argument makes
the resulting relation submodule constant and forces an
exponential-polynomial identity.  A common finite tree and a Cauchy
transform show that its combined inverse-branch densities vanish edge by
edge.  Finally, a simultaneous finite Galois cover realizes every such
functional relation by functoriality in the formal period space.
\end{abstract}

\medskip
\noindent\textit{2020 Mathematics Subject Classification.}
Primary 11J91; Secondary 11J81, 14F10, 30E20, 34M35.

\noindent\textit{Keywords.}
Exponential periods, formal periods, Kontsevich--Zagier period conjecture,
$E$-functions, affine line, twisted de Rham cohomology, Cauchy transforms.

\section{Introduction and main results}\label{sec:introduction}

Kontsevich and Zagier formulated the period conjecture as the assertion that
all identities between integral presentations of ordinary periods follow
from additivity, algebraic change of variables, and Stokes' formula
\cite{KontsevichZagier}.  In cohomological language this is the injectivity
of a formal period map; see
\cite[Chapters~12--13]{HuberMullerStach} and
\cite{CressonViuSos}.  Exponential periods replace an algebraic differential
form $\omega$ by $e^f\omega$, where $f$ is a regular function.  The de Rham
differential becomes
\[
  d_f=d+df\wedge,
\]
and Stokes' formula takes the twisted form
\[
  d(Ge^f)=e^f(dG+G\,df).
\]
General definitions and comparison results for exponential periods are
developed in \cite{CommelinHabeggerHuber}; the de Rham and rapid-decay
realizations for irregular connections on curves go back to
\cite{BlochEsnault,HienFlat}.  We use the terminology \emph{exponential
period function} in the sense of Fres\'an and Jossen
\cite{FresanJossen}.

This paper proves the linear period principle for compact exponential
periods on the affine line.  Let $q,r\in\Qbar[x]$, and let
\[
  \Delta=\sum_{\alpha\in S}d_\alpha[\alpha],
  \qquad S\subset\Qbar\text{ finite},\quad
  d_\alpha\in\Qbar,\quad \sum_\alpha d_\alpha=0.
\]
Choose a finite piecewise $C^1$ one-chain $\Gamma$ in $\C$ with
$\partial\Gamma=\Delta$, and put
\begin{equation}\label{eq:I}
  I_{q,r,\Delta}(z)=\int_\Gamma r(x)e^{zq(x)}\,dx.
\end{equation}
The one-form in \eqref{eq:I} has an entire primitive in $x$, so the function
is independent of the chosen chain and is entire in $z$.  The number of
interest is $I_{q,r,\Delta}(1)$; a nonzero algebraic evaluation parameter is
already absorbed into the polynomial $q$.

The formal period space used below is defined precisely in
\cref{sec:formal-space}.  It is generated by compact relative periods on
smooth affine curves and zero-dimensional exponential periods, modulo bilinearity, functoriality,
and twisted Stokes.  Let
$\Pcal_{\Aone}^{\mathrm{KZ,cpt}}$ denote its subspace generated by the
symbols attached to \eqref{eq:I} and by $e^\beta$, $\beta\in\Qbar$.

\begin{theorem}[Linear Kontsevich--Zagier theorem]\label{thm:main}
The period map
\[
  \per:\Pcal_{\Aone}^{\mathrm{KZ,cpt}}\longrightarrow\C
\]
is injective.  Thus every $\Qbar$-linear relation among compact exponential
periods on the affine line and zero-dimensional exponential periods is generated by bilinearity,
functoriality through smooth affine curves, and twisted Stokes.
\end{theorem}

The arithmetic core is a transfer theorem for a finite family of different
polynomials.  The degree condition on $s_j$ is the canonical twisted de
Rham normal form established in \cref{sec:reduction}.

\begin{theorem}[Transfer and separation of linear relations]
\label{thm:transfer}
For $1\le j\le N$, let $q_j,s_j\in\Qbar[x]$ and let $\Delta_j$ be a finite
degree-zero zero-cycle on $\Aone(\Qbar)$ with coefficients in $\Qbar$.
Assume
\[
  n_j:=\deg q_j\ge2,
  \qquad \deg s_j\le n_j-2,
\]
and put
\[
  F_j(z)=I_{q_j,s_j,\Delta_j}(z).
\]
Let $\beta_1,\ldots,\beta_M\in\Qbar$ be distinct.  For
$a_1,\ldots,a_N,b_1,\ldots,b_M\in\Qbar$, one has
\begin{equation}\label{eq:numerical-relation}
  \sum_{j=1}^{N}a_jF_j(1)
  +\sum_{\nu=1}^{M}b_\nu e^{\beta_\nu}=0
\end{equation}
if and only if
\begin{equation}\label{eq:functional-relation}
  \sum_{j=1}^{N}a_jF_j(z)\equiv0
\end{equation}
and
\[
  b_1=\cdots=b_M=0.
\]
\end{theorem}

The proof also gives a finite geometric criterion.  A common tree in the
targets of the $q_j$ carries inverse-branch densities $\psi_{j,e}$, defined
in \eqref{eq:individual-density}.  For any fixed common tree representation,
\cref{thm:transfer} strengthens to
\begin{equation}\label{eq:geometric-summary}
  \eqref{eq:numerical-relation}
  \quad\Longleftrightarrow\quad
  b_1=\cdots=b_M=0
  \ \text{and}\ 
  \sum_j a_j\psi_{j,e}\equiv0
  \text{ for every open edge }e.
\end{equation}
The vanishing density identities are then shown to be formal consequences
of functoriality on one simultaneous algebraic cover.  This is the step that
turns the $E$-function relation theorem into a Kontsevich--Zagier statement.

The companion paper \cite{LongEIT} treats $r=1$ for one fixed polynomial and
proves algebraic-value rigidity together with an exact description of
constant-coefficient functional relations.  The present paper adds arbitrary
polynomial differentials, different polynomials in one relation,
zero-dimensional exponential periods, and the formal-period interpretation.  For any fixed finite tuple
of $E$-functions, Fischler and Rivoal give an algorithm for the ideals of
functional and value relations \cite{FischlerRivoalEffective}; algorithms
for exceptional algebraic values and the required minimal differential
equations were developed by Adamczewski--Rivoal and
Bostan--Rivoal--Salvy \cite{AdamczewskiRivoal,BostanRivoalSalvy}.  The result
here is instead uniform in the polynomial and boundary data and identifies
the relations by a finite geometric criterion and by formal period
operations.  Delaygue's linear-independence theorem for $E$-values assumes
a separation condition on associated finite singularity sets
\cite{Delaygue}; no such separation is imposed here.

The conjunction of compact chains with the affine-line polynomial setting
is essential to the present argument.  In this setting the Taylor
coefficients of \eqref{eq:I} are algebraic and therefore lie within the
theory of $E$-functions.  General curves introduce ordinary period coefficients,
products lead to higher-dimensional periods, and noncompact rapid-decay
chains include Gamma values.  The precise boundary of the result is
recorded in \cref{sec:scope}.

The paper is organized as follows.  \Cref{sec:formal-space} defines the
formal period space.  \Cref{sec:reduction} gives the canonical reduction of
polynomial differentials.  \Cref{sec:moment-package,sec:direct-sums} record
the input from Paper~I and prove direct-sum rigidity.
\Cref{sec:transfer-proof} proves \cref{thm:transfer} by Beukers lifting and a
common-tree Cauchy transform.  \Cref{sec:formal-relations} realizes the
resulting functional relations formally on a simultaneous Galois cover.
The main theorem and applications follow in \cref{sec:main-proof}.

\section{The formal period space}\label{sec:formal-space}

For a smooth connected affine curve $Y/\Qbar$, put
\begin{equation}\label{eq:relative-homology}
  H_1^{\mathrm{rel}}(Y)
  =\varinjlim_{S\subset Y(\Qbar)\ \mathrm{finite}}
    H_1\bigl(Y(\C),S;\Qbar\bigr).
\end{equation}
The transition maps are induced by inclusions of finite sets.  Every class
is represented by a finite piecewise $C^1$ one-chain with coefficients in
$\Qbar$ and boundary supported on algebraic points.  Its boundary is a
well-defined finite degree-zero zero-cycle on $Y(\Qbar)$ with coefficients
in $\Qbar$.  A morphism $\rho:Y\to X$ induces a pushforward
$\rho_*:H_1^{\mathrm{rel}}(Y)\to H_1^{\mathrm{rel}}(X)$.

\begin{definition}[Compact curve-level formal exponential periods]
\label{def:curve-space}
Let $\Pcal_{\mathrm{curve}}^{\mathrm{KZ,cpt}}$ be the $\Qbar$-vector space
generated by symbols
\[
  \langle Y,f,\omega,\gamma\rangle,
\]
where $Y/\Qbar$ is a smooth connected affine curve,
$f\in\Ocal(Y)$, $\omega\in\Gamma(Y,\Omega^1_{Y/\Qbar})$, and
$\gamma\in H_1^{\mathrm{rel}}(Y)$, together with zero-dimensional symbols
$\mathsf E(\beta)$ for $\beta\in\Qbar$, modulo the following relations.

\begin{enumerate}[label=\textup{(K\arabic*)},leftmargin=3.2em]
\item \emph{Bilinearity.}  The symbol is $\Qbar$-linear in $\omega$ and
in $\gamma$.

\item \emph{Functoriality.}  If $\rho:Y\to X$ is a morphism of smooth
connected affine curves, $f\in\Ocal(X)$,
$\omega\in\Gamma(X,\Omega^1_{X/\Qbar})$, and
$\gamma\in H_1^{\mathrm{rel}}(Y)$, then
\begin{equation}\label{eq:functoriality}
  \langle Y,f\circ\rho,\rho^*\omega,\gamma\rangle
  =\langle X,f,\omega,\rho_*\gamma\rangle.
\end{equation}

\item \emph{Twisted Stokes.}  If $G\in\Ocal(Y)$ and
$\partial\gamma=\sum_p n_p[p]$, then
\begin{equation}\label{eq:twisted-stokes}
  \langle Y,f,dG+G\,df,\gamma\rangle
  =\sum_p n_pG(p)\mathsf E(f(p)).
\end{equation}
\end{enumerate}
\end{definition}

The assignments
\begin{equation}\label{eq:period-map}
  \per\bigl(\langle Y,f,\omega,\gamma\rangle\bigr)
  =\int_\gamma e^f\omega,
  \qquad
  \per(\mathsf E(\beta))=e^\beta
\end{equation}
define a well-defined $\Qbar$-linear period map.  Indeed, $e^f\omega$ is a
closed holomorphic one-form, \eqref{eq:functoriality} is change of variables,
and \eqref{eq:twisted-stokes} is the identity
$d(Ge^f)=e^f(dG+G\,df)$.

Let $\Delta$ be a finite degree-zero zero-cycle on $\Aone(\Qbar)$ with
coefficients in $\Qbar$.  The long exact sequence of the pair
$(\Aone(\C),\supp\Delta)$ shows that there is a unique class
\[
  \gamma_\Delta\in H_1^{\mathrm{rel}}(\Aone)
  \qquad\text{with}\qquad
  \partial\gamma_\Delta=\Delta,
\]
because $H_1(\Aone(\C);\Qbar)=0$ and $\Aone(\C)$ is connected.

\begin{definition}[Affine-line subspace]\label{def:affine-line-space}
Let $\Pcal_{\Aone}^{\mathrm{KZ,cpt}}$ be the subspace of
$\Pcal_{\mathrm{curve}}^{\mathrm{KZ,cpt}}$ generated by
\[
  \langle\Aone,q,r(x)\,dx,\gamma_\Delta\rangle,
  \qquad q,r\in\Qbar[x],
\]
for finite degree-zero zero-cycles $\Delta$, and by the zero-dimensional symbols
$\mathsf E(\beta)$.  We abbreviate
\begin{equation}\label{eq:P-symbol}
  \mathsf P(q,r,\Delta)
  =\langle\Aone,q,r(x)\,dx,\gamma_\Delta\rangle.
\end{equation}
Then
\[
  \per\bigl(\mathsf P(q,r,\Delta)\bigr)
  =I_{q,r,\Delta}(1).
\]
\end{definition}

\section{Twisted reduction of polynomial differentials}
\label{sec:reduction}

The first step is the one-variable twisted de Rham normal form.  It is
performed for the numerical period at $z=1$.

\begin{lemma}[Twisted polynomial reduction]\label{lem:twisted-reduction}
Let $q\in\Qbar[x]$ have degree $n\ge2$.  For every $r\in\Qbar[x]$ there
are unique polynomials $R,s\in\Qbar[x]$ such that
\begin{equation}\label{eq:twisted-reduction}
  r=R'+q'R+s,
  \qquad \deg s\le n-2.
\end{equation}
For every finite degree-zero zero-cycle
$\Delta=\sum_\alpha d_\alpha[\alpha]$, one has in
$\Pcal_{\Aone}^{\mathrm{KZ,cpt}}$
\begin{equation}\label{eq:formal-reduction}
  \mathsf P(q,r,\Delta)
  =\mathsf P(q,s,\Delta)
   +\sum_\alpha d_\alpha R(\alpha)\mathsf E(q(\alpha)).
\end{equation}
\end{lemma}

\begin{proof}
Write $q'(x)=cx^{n-1}+O(x^{n-2})$ with $c\ne0$.  If $r$ has degree
$m\ge n-1$ and leading term $ax^m$, then
\[
  R_0=\frac ac x^{m-n+1}
\]
has the property that $R_0'+q'R_0$ has leading term $ax^m$.  Subtracting
this polynomial lowers the degree, and iteration proves existence.

For uniqueness, suppose $R'+q'R+s=0$ with $\deg s\le n-2$.  If $R\ne0$,
then
\[
  \deg(R'+q'R)=\deg R+n-1\ge n-1,
\]
since the leading term of $q'R$ cannot cancel with $R'$.  Hence $R=0$ and
then $s=0$.

Finally, apply twisted Stokes on $\Aone$ to
\[
  d(Re^q)=e^q(R'+q'R)\,dx.
\]
This gives \eqref{eq:formal-reduction}.
\end{proof}

\begin{lemma}[Polynomials of degree at most one]\label{lem:small-degree}
Let $q,r\in\Qbar[x]$ and let $\Delta=\sum_\alpha d_\alpha[\alpha]$.
\begin{enumerate}[label=\textup{(\roman*)},leftmargin=2.4em]
\item If $q=\beta$ is constant, then
\[
  \mathsf P(q,r,\Delta)
  =\left(\sum_\alpha d_\alpha R(\alpha)\right)\mathsf E(\beta),
\]
where $R'=r$.
\item If $q(x)=ax+b$ with $a\ne0$, there is a unique
$R\in\Qbar[x]$ satisfying $r=R'+aR$, and
\begin{equation}\label{eq:linear-reduction}
  \mathsf P(q,r,\Delta)
  =\sum_\alpha d_\alpha R(\alpha)\mathsf E(q(\alpha)).
\end{equation}
\end{enumerate}
\end{lemma}

\begin{proof}
The constant case follows from twisted Stokes with a polynomial primitive
of $r$.  In the linear case, the map
\[
  \Qbar[x]\longrightarrow\Qbar[x],
  \qquad R\longmapsto R'+aR,
\]
preserves the degree filtration and induces multiplication by $a$ on every
associated graded piece.  It is therefore an automorphism.  Apply twisted
Stokes once more.
\end{proof}

\begin{corollary}[Formal normal form]\label{cor:normal-form}
Every element of $\Pcal_{\Aone}^{\mathrm{KZ,cpt}}$ is represented by a
finite $\Qbar$-linear combination of symbols
\[
  \mathsf P(q,s,\Delta),
  \qquad \deg q\ge2,
  \qquad \deg s\le\deg q-2,
\]
and zero-dimensional symbols $\mathsf E(\beta)$.
\end{corollary}

\begin{remark}\label{rem:reduction-at-one}
The reduction \eqref{eq:twisted-reduction} is canonical for the period at
$z=1$.  For the deformation \eqref{eq:I}, one has
\[
  d(Re^{zq})=e^{zq}(R'+zq'R)\,dx,
\]
so an unreduced differential does not retain the same decomposition as $z$
varies.  This is why \cref{thm:transfer} is stated for reduced
polynomials $s_j$, while the general formal theorem first applies
\cref{lem:twisted-reduction} at $z=1$.

For example, take $q(x)=x(x-1)$, $r(x)=2x^2-x+1$, and
$\Delta=[1]-[0]$.  Then $r=R'+q'R$ with $R=x$, so $s=0$ and
\[
  \int_0^1 r(x)e^{q(x)}\,dx
  =\bigl[xe^{q(x)}\bigr]_0^1=1.
\]
Nevertheless, $I_{q,r,\Delta}(0)=7/6$, so the deformed function is not
identically zero.  The algebraic value at $1$ is entirely an endpoint term;
the reduced transfer theorem cannot be applied to the unreduced function.
\end{remark}

\section{The moment package from Paper I}\label{sec:moment-package}

We recall only the properties of the single-polynomial moment system needed
below.  An entire function
$F(z)=\sum_{m\ge0}a_mz^m/m!$ is an $E$-function if its coefficients are
algebraic, their joint heights grow at most linearly, and $F$ satisfies a
nonzero linear differential equation over $\Qbar[z]$.

Let $q\in\Qbar[x]$ have degree $n\ge2$, let $\Delta$ be a finite
degree-zero zero-cycle, and choose a chain $\Gamma$ with
$\partial\Gamma=\Delta$.  Put
\begin{equation}\label{eq:moment-vector}
  U_j(z)=\int_\Gamma x^j e^{zq(x)}\,dx,
  \qquad 0\le j\le n-2,
  \qquad
  \mathbf U=(U_0,\ldots,U_{n-2})^{\mathsf T}.
\end{equation}

\begin{theorem}[Moment package from Paper I]\label{thm:moment-package}
There are matrices $C_q,D_q\in M_{n-1}(\Qbar)$, a finite set
\[
  \Lambda=\{0\}\cup q(\supp\Delta),
\]
and vectors $\mathbf b_\lambda\in\Qbar^{n-1}$ such that
\begin{equation}\label{eq:moment-system}
  \mathbf U'
  =\left(C_q+\frac{D_q}{z}\right)\mathbf U
   +\frac1z\sum_{\lambda\in\Lambda}
      \mathbf b_\lambda e^{\lambda z}.
\end{equation}
Moreover:
\begin{enumerate}[label=\textup{(\roman*)},leftmargin=2.4em]
\item every coordinate $U_j$ is an $E$-function;
\item the homogeneous differential module
\[
  \mathbf Y'=\left(C_q+\frac{D_q}{z}\right)\mathbf Y
\]
is semisimple over $\C(z)$;
\item every eigenvalue of $D_q$ lies strictly between $-1$ and $0$;
\item $C_q$ is semisimple, and for every eigenvalue $\lambda$ of $C_q$,
every eigenvalue of the compression
\[
  \pi_\lambda D_q|_{\ker(C_q-\lambda I)}
\]
lies strictly between $-1$ and $0$, where $\pi_\lambda$ is the spectral
projection of $C_q$.
\end{enumerate}
\end{theorem}

\begin{proof}
This is the unnormalized moment system of
\cite[Proposition~2.2]{LongEIT}, together with semisimplicity from
\cite[Proposition~3.1]{LongEIT}.  The global residue spectrum is computed in
\cite[Lemma~2.1]{LongEIT}; the semisimplicity of $C_q$ and the compressed
local spectra are established in the proof of
\cite[Proposition~3.2]{LongEIT}.  In particular, no affine normalization of
$q$ is required.
\end{proof}

We shall also use the following standard results.

\begin{theorem}[Beukers, linear lifting]\label{thm:Beukers}
Let $F_1,\ldots,F_s$ be $E$-functions forming a vector solution of
\[
  \mathbf F'=A(z)\mathbf F,
  \qquad A(z)\in M_s(\Qbar(z)),
\]
and let $T(z)\in\Qbar[z]$ be a common denominator of the entries of $A$.
If $\xi\in\Qbar$ satisfies $\xi T(\xi)\ne0$ and
\[
  \alpha_1F_1(\xi)+\cdots+\alpha_sF_s(\xi)=0,
  \qquad \alpha_i\in\Qbar,
\]
then there are polynomials $A_1(z),\ldots,A_s(z)\in\Qbar[z]$ such that
\[
  \sum_{i=1}^s A_i(z)F_i(z)\equiv0,
  \qquad A_i(\xi)=\alpha_i.
\]
\end{theorem}

\begin{proof}
This is the linear lifting theorem
\cite[Theorem~3.2 and Lemma~3.1]{Beukers}.  No linear independence
hypothesis on the coordinates is required.
\end{proof}

\begin{lemma}[Distinct exponential functions]\label{lem:distinct-exponentials}
If $\lambda_1,\ldots,\lambda_s\in\C$ are distinct, then
$e^{\lambda_1z},\ldots,e^{\lambda_sz}$ are linearly independent over
$\C(z)$.
\end{lemma}

\begin{proof}
Choose a nontrivial relation with the fewest nonzero terms.  Differentiate
and eliminate one term.  Minimality forces the logarithmic derivative of a
nonzero rational function to be a nonzero constant, which is impossible at
infinity.
\end{proof}

\begin{theorem}[Lindemann--Weierstrass, linear form]
\label{thm:Lindemann-Weierstrass}
If $\beta_1,\ldots,\beta_s\in\Qbar$ are distinct, then
$e^{\beta_1},\ldots,e^{\beta_s}$ are linearly independent over $\Qbar$.
\end{theorem}

\begin{proof}
See, for example, \cite[Chapter~1]{Baker}.
\end{proof}

\section{Direct sums and constant relation submodules}
\label{sec:direct-sums}

Let $q_1,\ldots,q_N\in\Qbar[x]$ have degrees $n_j\ge2$.  For each $j$,
let $C_j,D_j$ be the matrices supplied by \cref{thm:moment-package}, and
put
\begin{equation}\label{eq:block-CD}
  C=\bigoplus_{j=1}^{N}C_j,
  \qquad
  D=\bigoplus_{j=1}^{N}D_j,
  \qquad
  \mu=\sum_{j=1}^{N}(n_j-1).
\end{equation}
The direct sum of the homogeneous moment modules is semisimple.  The next
result shows that its rational horizontal endomorphisms are still constant,
even when different polynomials have common critical values.

\begin{theorem}[Direct-sum endomorphism rigidity]
\label{thm:direct-sum-rigidity}
If $P(z)\in M_\mu(\C(z))$ satisfies
\begin{equation}\label{eq:endomorphism-equation}
  P'=\left[C+\frac Dz,P\right],
\end{equation}
then $P$ is constant and commutes with both $C$ and $D$.
\end{theorem}

\begin{proof}
A pole of order $m\ge1$ at a point $z_0\ne0$ is impossible: differentiation
raises its order to $m+1$, whereas the right side of
\eqref{eq:endomorphism-equation} has order at most $m$.

Suppose that $P$ has a pole of order $m\ge1$ at $0$ and write
\[
  P=z^{-m}P_{-m}+O(z^{-m+1}),
  \qquad P_{-m}\ne0.
\]
The coefficient of $z^{-m-1}$ gives
\begin{equation}\label{eq:zero-resonance}
  -mP_{-m}=[D,P_{-m}].
\end{equation}
Every eigenvalue of $D$ lies in $(-1,0)$ by
\cref{thm:moment-package}(iii).  Hence every eigenvalue of
$X\mapsto[D,X]$ is a difference of two numbers in $(-1,0)$ and lies in
$(-1,1)$.  It cannot equal the negative integer $-m$.  Thus $P$ has no
finite pole and is a polynomial matrix.

Assume that
\[
  P(z)=P_dz^d+P_{d-1}z^{d-1}+\cdots,
  \qquad d\ge1,
  \quad P_d\ne0.
\]
Comparing the coefficients of $z^d$ and $z^{d-1}$ in
\eqref{eq:endomorphism-equation} gives
\begin{align}
  [C,P_d]&=0,\label{eq:leading-C}\\
  dP_d&=[C,P_{d-1}]+[D,P_d].\label{eq:next-leading}
\end{align}
The matrix $C$ is semisimple.  For an eigenvalue $\lambda$ of $C$, let
$\pi_\lambda$ be its spectral projection and put
\[
  P_{d,\lambda}=\pi_\lambda P_d\pi_\lambda,
  \qquad
  \overline D_\lambda
  =\pi_\lambda D|_{\ker(C-\lambda I)}.
\]
Equation \eqref{eq:leading-C} makes $P_d$ block diagonal for the
$C$-eigenspace decomposition.  Compressing \eqref{eq:next-leading} to the
$\lambda$-eigenspace eliminates the $C$-commutator and yields
\begin{equation}\label{eq:infinity-resonance}
  dP_{d,\lambda}
  =[\overline D_\lambda,P_{d,\lambda}].
\end{equation}
For at least one $\lambda$, the block $P_{d,\lambda}$ is nonzero.

The matrix $\overline D_\lambda$ is the direct sum of the compressed blocks
belonging to all critical points of all $q_j$ with critical value
$\lambda$.  By \cref{thm:moment-package}(iv), all of its eigenvalues lie in
$(-1,0)$.  Therefore every eigenvalue of
$X\mapsto[\overline D_\lambda,X]$ lies in $(-1,1)$, contradicting
\eqref{eq:infinity-resonance} because $d$ is a positive integer.  Hence
$d=0$.

Thus $P$ is constant.  Substitution in
\eqref{eq:endomorphism-equation} gives $[C,P]=[D,P]=0$.
\end{proof}

\begin{corollary}[Constant relation submodules]
\label{cor:constant-submodules}
Let $\Mcal$ be the row differential module with connection
\[
  \boldsymbol\phi\longmapsto
  \boldsymbol\phi'
  +\boldsymbol\phi\left(C+\frac Dz\right).
\]
Every differential submodule $R\subseteq\Mcal$ has the form
\[
  R=W\otimes_\C\C(z)
\]
for a constant subspace $W\subseteq\C^{1\times\mu}$ invariant under right
multiplication by $C$ and $D$.
\end{corollary}

\begin{proof}
The direct-sum module and its dual are semisimple.  Hence $R$ has a
horizontal complement and a rational horizontal idempotent projector.  In
row coordinates the projector acts by right multiplication by a matrix
$P(z)$ satisfying \eqref{eq:endomorphism-equation}.  By
\cref{thm:direct-sum-rigidity}, $P$ is constant and commutes with $C,D$.
Take $W$ to be its constant row image.
\end{proof}

\section{Transfer of linear relations}
\label{sec:transfer-proof}

Fix data as in \cref{thm:transfer}.  For each $j$, choose a chain with
boundary $\Delta_j$ and let $\mathbf U_j$ be its moment vector
\eqref{eq:moment-vector}.  Put
\[
  \mathbf U=\bigoplus_{j=1}^{N}\mathbf U_j.
\]
Since $\deg s_j\le n_j-2$, there is a constant row
$v\in\Qbar^{1\times\mu}$ such that
\begin{equation}\label{eq:F-vU}
  F(z):=\sum_{j=1}^{N}a_jF_j(z)=v\mathbf U(z).
\end{equation}
Let $\Lambda\subset\Qbar$ be the finite set consisting of $0$, all values
$q_j(\alpha)$ with $\alpha\in\supp\Delta_j$, and all $\beta_\nu$.  Put
\[
  \Ecal_\Lambda
  =\Span_{\C(z)}\{e^{\lambda z}:\lambda\in\Lambda\}.
\]
The direct-sum inhomogeneous system has the form
\begin{equation}\label{eq:direct-inhomogeneous-system}
  \mathbf U'
  =\left(C+\frac Dz\right)\mathbf U
   +\frac1z\sum_{\lambda\in\Lambda}
      \boldsymbol\eta_\lambda e^{\lambda z},
\end{equation}

\begin{proposition}[From a numerical relation to an exponential polynomial]
\label{prop:value-to-exp-poly}
If \eqref{eq:numerical-relation} holds, then
\begin{equation}\label{eq:F-exp-polynomial}
  F(z)=\sum_{\lambda\in\Lambda}p_\lambda(z)e^{\lambda z}
\end{equation}
for polynomials $p_\lambda\in\Qbar[z]$.
\end{proposition}

\begin{proof}
Adjoin the functions $e^{\lambda z}$, $\lambda\in\Lambda$, to
\eqref{eq:direct-inhomogeneous-system}.  The resulting vector consists of
$E$-functions and satisfies a homogeneous first-order system over
$\Qbar(z)$ with common denominator $z$.  By \cref{thm:Beukers}, the
relation at $z=1$ lifts to
\begin{equation}\label{eq:lifted-relation}
  \boldsymbol\phi(z)\mathbf U(z)
  +\sum_{\lambda\in\Lambda}c_\lambda(z)e^{\lambda z}=0,
  \qquad
  \boldsymbol\phi(1)=v,
\end{equation}
where $\boldsymbol\phi$ and the $c_\lambda$ have coefficients in
$\Qbar[z]$.

Define
\begin{equation}\label{eq:relation-module}
  \Rcal=
  \left\{
    \boldsymbol\psi\in\C(z)^{1\times\mu}:
    \boldsymbol\psi\mathbf U\in\Ecal_\Lambda
  \right\}.
\end{equation}
This is a differential submodule of the dual homogeneous moment module.
Indeed, if $A=C+D/z$ and
$\mathbf f=z^{-1}\sum_\lambda\boldsymbol\eta_\lambda e^{\lambda z}$,
then
\[
  (\boldsymbol\psi\mathbf U)'
  =(\boldsymbol\psi'+\boldsymbol\psi A)\mathbf U
   +\boldsymbol\psi\mathbf f.
\]
The left side and the last term belong to $\Ecal_\Lambda$, so
$\boldsymbol\psi'+\boldsymbol\psi A\in\Rcal$.

By \cref{cor:constant-submodules},
\[
  \Rcal=W\otimes_\C\C(z)
\]
for a constant $C,D$-invariant row subspace $W$.  The row
$\boldsymbol\phi$ in \eqref{eq:lifted-relation} belongs to $\Rcal$ and is
regular at $1$.  Choose a constant matrix $T_W$ whose columns span the
annihilator of $W$.  Then $\boldsymbol\phi(z)T_W=0$, and evaluation at
$1$ gives $vT_W=0$.  Thus $v\in W$.  Hence the constant row
$v$ belongs to $\Rcal$, and
\begin{equation}\label{eq:F-in-E}
  F=v\mathbf U\in\Ecal_\Lambda.
\end{equation}

Choose a constant full-row-rank matrix $S$ whose rows form a basis of $W$,
and put $\mathbf V=S\mathbf U$.  The invariance of $W$ gives constant
matrices $C_W,D_W$ satisfying
\[
  SC=C_WS,
  \qquad
  SD=D_WS.
\]
Thus
\begin{equation}\label{eq:restricted-system}
  \mathbf V'
  =\left(C_W+\frac{D_W}{z}\right)\mathbf V
   +\frac1z\sum_{\lambda\in\Lambda}
       \mathbf b_{W,\lambda}e^{\lambda z}.
\end{equation}
Every component of $\mathbf V$ lies in $\Ecal_\Lambda$.  By
\cref{lem:distinct-exponentials}, there are unique rational vectors
$\mathbf r_\lambda(z)$ such that
\begin{equation}\label{eq:V-expansion}
  \mathbf V(z)=\sum_{\lambda\in\Lambda}
    \mathbf r_\lambda(z)e^{\lambda z}.
\end{equation}
Comparison of the independent exponential coefficients in
\eqref{eq:restricted-system} gives
\begin{equation}\label{eq:r-equation}
  \mathbf r_\lambda'
  =(C_W-\lambda I)\mathbf r_\lambda
   +\frac{D_W}{z}\mathbf r_\lambda
   +\frac{\mathbf b_{W,\lambda}}{z}.
\end{equation}

A pole of $\mathbf r_\lambda$ at $z_0\ne0$ is impossible because
its derivative has one order more than every term on the right.  If
\[
  \mathbf r_\lambda=z^{-m}u+O(z^{-m+1}),
  \qquad m\ge1,
  \quad u\ne0,
\]
then the coefficient of $z^{-m-1}$ in \eqref{eq:r-equation} gives
$-mu=D_Wu$.  The eigenvalues of $D_W$ are among those of $D$ and lie in
$(-1,0)$, so this is impossible.  Thus every $\mathbf r_\lambda$ is a
polynomial vector, and \eqref{eq:F-in-E} gives
\eqref{eq:F-exp-polynomial} with $p_\lambda\in\C[z]$.

Finally, the coefficients of the $p_\lambda$ are algebraic by a finite
Hermite-interpolation argument.  If every $p_\lambda$ is zero, there is
nothing to prove.  Otherwise choose an integer $\ell\ge1$ such that
$\deg p_\lambda<\ell$ for every $\lambda$, write
\[
  p_\lambda(z)=\sum_{k=0}^{\ell-1}c_{\lambda,k}z^k,
  \qquad h=\ell|\Lambda|,
\]
and differentiate \eqref{eq:F-exp-polynomial} at zero.  For $0\le m<h$,
\[
  F^{(m)}(0)
  =\sum_{\lambda\in\Lambda}\sum_{k=0}^{\ell-1}
    c_{\lambda,k}
    \left.\frac{d^k}{dt^k}t^m\right|_{t=\lambda}.
\]
The coefficient matrix is the transpose of the matrix of the Hermite
jet-evaluation map
\[
  \Qbar[t]_{<h}\longrightarrow
  \bigoplus_{\lambda\in\Lambda}\Qbar^\ell,
  \qquad
  H\longmapsto
  \bigl(H^{(k)}(\lambda)\bigr)_{\lambda,\,0\le k<\ell}.
\]
This map is injective: an element of its kernel is divisible by
$\prod_{\lambda\in\Lambda}(t-\lambda)^\ell$, whose degree is $h$.
Since the source and target both have dimension $h$, the matrix is
invertible over $\Qbar$.  Every $F^{(m)}(0)$ is algebraic, so the unique
solution has $c_{\lambda,k}\in\Qbar$.  Hence
$p_\lambda\in\Qbar[z]$.
\end{proof}

\section{A common tree and the Cauchy-transform criterion}
\label{sec:common-tree}

We next show that a reduced linear combination of the functions under
consideration cannot be a nonzero exponential polynomial.

\begin{lemma}[Connectivity of a polynomial branch tree]
\label{lem:branch-tree-connectivity}
Let $q\in\C[x]$ be nonconstant, let $\Sigma$ be its set of finite critical
values, and choose $c\in\C\setminus\Sigma$.  For each
$\sigma\in\Sigma$, choose a simple arc $A_\sigma$ from $c$ to $\sigma$
whose interior avoids $\Sigma$, with pairwise disjoint interiors.  Put
\[
  T_\Sigma=\bigcup_{\sigma\in\Sigma}A_\sigma,
\]
with $T_\Sigma=\{c\}$ when $\Sigma$ is empty.  Then
$q^{-1}(T_\Sigma)$ is connected.
\end{lemma}

\begin{proof}
If $\Sigma$ is empty, then $q$ is linear.  Otherwise let
$E=q^{-1}(c)$.  For each $\sigma$, a standard loop based at $c$ that runs
along $A_\sigma$, circles $\sigma$, and returns along the other side has a
monodromy permutation $g_\sigma$ of $E$.  The local model
$u\mapsto u^r$ shows that the lifts of $A_\sigma$ whose initial sheets
belong to one cycle of $g_\sigma$ meet at the corresponding point above
$\sigma$.  It follows that the connected components of
$q^{-1}(T_\Sigma)$ meet $E$ in the orbits of the group generated by the
$g_\sigma$.

The loops just described generate $\pi_1(\C\setminus\Sigma,c)$, and
$q^{-1}(\C\setminus\Sigma)=\C\setminus q^{-1}(\Sigma)$ is connected.
Thus their monodromy action on $E$ is transitive.  Hence all points of $E$
lie in one component of $q^{-1}(T_\Sigma)$, and every other point is joined
to $E$ by a lifted subarc.
\end{proof}

\begin{lemma}[Common tree representation]
\label{lem:common-tree}
Let $q_1,\ldots,q_N\in\Qbar[x]$ be nonconstant, and let $\Delta_j$ be
finite degree-zero zero-cycles.  There is a finite embedded polygonal tree
$T\subset\C$, with algebraic vertices, such that:
\begin{enumerate}[label=\textup{(\roman*)},leftmargin=2.4em]
\item $T$ contains every finite critical value of every $q_j$ and every
value $q_j(\alpha)$ with $\alpha\in\supp\Delta_j$;
\item for each $j$ there is a finite one-chain $\Gamma_j$ in
$q_j^{-1}(T)$, with coefficients in $\Qbar$, such that
$\partial\Gamma_j=\Delta_j$;
\item over every open edge of $T$, the chain $\Gamma_j$ is a finite
$\Qbar$-linear combination of analytic inverse branches of $q_j$.
\end{enumerate}
\end{lemma}

\begin{proof}
Choose an algebraic point $c$ that is neither a critical value nor an
endpoint value for any $q_j$.  Join $c$ to all distinct critical and
endpoint values by simple polygonal arcs with pairwise disjoint interiors,
using algebraic auxiliary vertices.  Their union is a finite tree $T$.

For fixed $j$, the union of the arms ending at the critical values of
$q_j$ has connected inverse image by
\cref{lem:branch-tree-connectivity}.  Every other arm has regular interior,
and all its lifts attach above $c$; adjoining them preserves connectivity.
Thus $q_j^{-1}(T)$ is connected.

Choose $x_{j,*}\in q_j^{-1}(c)$ and write
$\Delta_j=\sum_\alpha d_{j,\alpha}[\alpha]$.  For each endpoint
$\alpha$, choose an edge path $P_{j,\alpha}$ in $q_j^{-1}(T)$ from
$x_{j,*}$ to $\alpha$, and put
\[
  \Gamma_j=\sum_\alpha d_{j,\alpha}P_{j,\alpha}.
\]
The coefficients at $x_{j,*}$ cancel because $\deg\Delta_j=0$, so
$\partial\Gamma_j=\Delta_j$.  Over the interior of any edge, all $q_j$ are
unramified, giving the asserted inverse-branch decomposition.
\end{proof}

Orient every open edge $e$ of $T$.  Write the inverse branches of $q_j$
over $e$ as
\[
  x_{j,e,1},\ldots,x_{j,e,n_j},
  \qquad q_j(x_{j,e,i}(t))=t,
\]
and let $f_{j,e,i}\in\Qbar$ be the coefficient with which $\Gamma_j$
traverses the corresponding lift.  Define
\begin{equation}\label{eq:individual-density}
  \psi_{j,e}(t)
  =\sum_i f_{j,e,i}
    \frac{s_j(x_{j,e,i}(t))}{q_j'(x_{j,e,i}(t))}
\end{equation}
and
\begin{equation}\label{eq:combined-density}
  \Psi_e(t)=\sum_{j=1}^{N}a_j\psi_{j,e}(t).
\end{equation}
These functions are analytic on the open edge and have integrable algebraic
singularities at its endpoints.  At a point of ramification index $r$, an
inverse-branch derivative is $O(|t-t_0|^{1/r-1})$.

\begin{theorem}[Combined moment rigidity]
\label{thm:combined-moment-rigidity}
Let $0\ne H\in\C[t]$.  If
\begin{equation}\label{eq:combined-moments}
  \sum_{j=1}^{N}a_j
  \int_{\Gamma_j}s_j(x)q_j(x)^mH(q_j(x))\,dx=0
  \qquad(m\ge0),
\end{equation}
then
\begin{equation}\label{eq:density-vanishing}
  \Psi_e\equiv0
  \qquad\text{for every open edge }e\subset T.
\end{equation}
Consequently,
\begin{equation}\label{eq:functional-vanishing-conclusion}
  \sum_{j=1}^{N}a_jI_{q_j,s_j,\Delta_j}(z)\equiv0.
\end{equation}
Conversely, \eqref{eq:functional-vanishing-conclusion} implies
\eqref{eq:density-vanishing}.
\end{theorem}

\begin{proof}
For $w\in\C\setminus T$, define
\begin{equation}\label{eq:combined-Cauchy-transform}
  \mathcal C_H(w)
  =\sum_{j=1}^{N}a_j
    \int_{\Gamma_j}
      \frac{s_j(x)H(q_j(x))}{w-q_j(x)}\,dx.
\end{equation}
For sufficiently large $|w|$, the geometric expansion of the denominator
converges uniformly on all chains.  Equation \eqref{eq:combined-moments}
therefore gives $\mathcal C_H(w)=0$.  The complement of a finite plane tree
is connected, so the identity theorem yields
\begin{equation}\label{eq:Cauchy-zero}
  \mathcal C_H\equiv0
  \qquad\text{on }\C\setminus T.
\end{equation}

Changing variables $t=q_j(x)$ on each lifted edge gives the absolutely
convergent representation
\begin{equation}\label{eq:Cauchy-edge-representation}
  \mathcal C_H(w)
  =\sum_e\int_e\frac{H(t)\Psi_e(t)}{w-t}\,dt.
\end{equation}
Fix an interior point $t_0$ of an edge $e$ with $H(t_0)\ne0$.  Choose a
small disk $B$ such that $B\cap T$ is a compact subarc
$\eta\Subset e$ containing $t_0$.  All relevant inverse branches extend
holomorphically across $B$, so
$\rho(t)=H(t)\Psi_e(t)$ is holomorphic there.  The contribution of
$T\setminus\eta$ extends holomorphically across $B$, while
\begin{equation}\label{eq:local-Cauchy-splitting}
  \int_\eta\frac{\rho(t)}{w-t}\,dt
  =\rho(w)\int_\eta\frac{dt}{w-t}
   -\int_\eta\frac{\rho(t)-\rho(w)}{t-w}\,dt.
\end{equation}
The second term extends holomorphically across $\eta$.  The two lateral
values of the first differ by
$2\pi i\varepsilon_e\rho(t_0)$, where
$\varepsilon_e\in\{\pm1\}$ depends on the orientation.  Both lateral
values of $\mathcal C_H$ vanish by \eqref{eq:Cauchy-zero}; hence
$H(t_0)\Psi_e(t_0)=0$.  Analyticity and the finiteness of the zero set of
$H$ give \eqref{eq:density-vanishing}.

Finally,
\[
  \sum_j a_jI_{q_j,s_j,\Delta_j}(z)
  =\sum_e\int_e e^{zt}\Psi_e(t)\,dt=0.
\]
Conversely, functional vanishing makes every Taylor coefficient at $z=0$
zero.  This is \eqref{eq:combined-moments} with $H=1$, and the preceding
argument again gives \eqref{eq:density-vanishing}.
\end{proof}

The proof is a common-tree version of the one-polynomial Cauchy-transform
criterion in \cite{LongEIT}; it is closely related to the polynomial moment
and algebraic Cauchy-transform methods of
\cite{PakovichMuzychuk,PakovichRoytvarfYomdin,GavrilovPakovich}.

\begin{corollary}[Exponential-polynomial exclusion]
\label{cor:exp-poly-exclusion}
If
\[
  F(z)=\sum_{j=1}^{N}a_jI_{q_j,s_j,\Delta_j}(z)
\]
is an exponential polynomial, then $F\equiv0$.
\end{corollary}

\begin{proof}
If $F\equiv0$, there is nothing to prove.  Otherwise write
\[
  F(z)=\sum_{\lambda\in\Lambda}p_\lambda(z)e^{\lambda z}
\]
with $\Lambda$ finite and at least one $p_\lambda$ nonzero.  Put
\[
  L=1+\max_{p_\lambda\ne0}\deg p_\lambda,
  \qquad
  \mathscr D=\prod_{\lambda\in\Lambda}(\partial_z-\lambda)^L,
  \qquad
  H(t)=\prod_{\lambda\in\Lambda}(t-\lambda)^L.
\]
Then $\mathscr DF=0$, and differentiation under the integral sign gives
\[
  0=\sum_j a_j\int_{\Gamma_j}
     s_j(x)H(q_j(x))e^{zq_j(x)}\,dx.
\]
Its Taylor coefficients at $z=0$ are precisely
\eqref{eq:combined-moments}.  Apply
\cref{thm:combined-moment-rigidity}.
\end{proof}

\begin{proof}[Proof of \cref{thm:transfer}]
Assume \eqref{eq:numerical-relation}.  By
\cref{prop:value-to-exp-poly}, the function
$F=\sum_j a_jF_j$ is an exponential polynomial.  By
\cref{cor:exp-poly-exclusion}, $F\equiv0$.  The numerical relation now
reduces to
\[
  \sum_{\nu=1}^{M}b_\nu e^{\beta_\nu}=0,
\]
and the distinctness of the $\beta_\nu$ together with
\cref{thm:Lindemann-Weierstrass} gives $b_1=\cdots=b_M=0$.
The converse follows by evaluating the functional relation at $z=1$.
\end{proof}

\begin{corollary}[Geometric criterion for numerical relations]
\label{cor:geometric-criterion}
For any common tree representation from \cref{lem:common-tree}, the
numerical relation \eqref{eq:numerical-relation} holds if and only if
\[
  b_1=\cdots=b_M=0
  \qquad\text{and}\qquad
  \sum_j a_j\psi_{j,e}\equiv0
  \quad\text{for every open edge }e.
\]
\end{corollary}

\begin{proof}
Combine \cref{thm:transfer} with both directions of
\cref{thm:combined-moment-rigidity}.
\end{proof}

\begin{corollary}[Linear independence after adjoining zero-dimensional exponential periods]
\label{cor:independence-with-points}
Let $F_j=I_{q_j,s_j,\Delta_j}$ satisfy the reduced hypotheses of
\cref{thm:transfer}, and let $\beta_1,\ldots,\beta_M\in\Qbar$ be
distinct.  If $F_1,\ldots,F_N$ are linearly independent over $\Qbar$ as
entire functions, then
\[
  F_1(1),\ldots,F_N(1),e^{\beta_1},\ldots,e^{\beta_M}
\]
are linearly independent over $\Qbar$.
\end{corollary}

\begin{proof}
Apply \cref{thm:transfer} to an arbitrary numerical relation.
\end{proof}

\begin{remark}[Algebraic evaluation points]
\label{rem:absorbed-parameters}
No extra evaluation parameters are needed in the statements.  For
$\xi\in\Qbar^\times$,
\[
  I_{q,s,\Delta}(\xi z)=I_{\xi q,s,\Delta}(z).
\]
Thus separate nonzero algebraic arguments for a finite family are absorbed
individually into the corresponding polynomials.
\end{remark}

\section{Functional relations in the formal period space}
\label{sec:formal-relations}

The density criterion is analytic.  We now show that every vanishing density
identity is already a consequence of functoriality and bilinearity in the
formal period space.

\begin{lemma}[Simultaneous finite Galois cover]
\label{lem:simultaneous-cover}
Let $q_1,\ldots,q_N\in\Qbar[x]$ be nonconstant, let $\Sigma$ be the union
of their finite critical-value sets, and put
$U=\Aone_t\setminus\Sigma$.  There exist a finite Galois extension
$L/\Qbar(t)$, a smooth connected affine curve $Y/\Qbar$, and a finite
morphism
\[
  \pi:Y\longrightarrow\Aone_t
\]
with the following property.  For every $j$ and every
$\Qbar(t)$-embedding
\[
  \sigma:\Qbar(x_j)\hookrightarrow L,
  \qquad t=q_j(x_j),
\]
there is a morphism
$\rho_{j,\sigma}:Y\to\Aone_x$ satisfying
\begin{equation}\label{eq:cover-factorization}
  q_j\circ\rho_{j,\sigma}=\pi.
\end{equation}
The morphism $\pi$ is finite \'{e}tale over $U$.
If $B\subset U(\C)$ is a nonempty, connected, simply connected open set,
and $\widetilde B$ is a connected component of $\pi^{-1}(B)$, then the
functions
\[
  \rho_{j,\sigma}\circ
  (\pi|_{\widetilde B})^{-1}
\]
are precisely the analytic inverse branches of $q_j$ on $B$.
\end{lemma}

\begin{proof}
For each $j$, put $K_j=\Qbar(x_j)$ with $t=q_j(x_j)$.  The extension
$K_j/\Qbar(t)$ is finite separable of degree $\deg q_j$.  Inside a fixed algebraic closure of $\Qbar(t)$, let
$K_j^{\mathrm{gal}}$ be the Galois closure of $K_j/\Qbar(t)$, and take
\[
  L=K_1^{\mathrm{gal}}\cdots K_N^{\mathrm{gal}}
\]
to be their compositum itself.  This is a finite Galois extension of
$\Qbar(t)$.  Let $Y$ be the normalization of $\Aone_t$ in $L$.  Since $L$ is a field, $Y$ is connected; normalization
is finite over the affine line, so $Y$ is affine; and a normal curve in
characteristic zero is smooth.  Let $\pi$ be the normalization morphism.

For every embedding $\sigma:K_j\hookrightarrow L$, the element
$\sigma(x_j)$ is integral over $\Qbar[t]$, because it satisfies the monic
equation obtained from $q_j(X)-t=0$ after division by the leading
coefficient.  Hence $\sigma(x_j)$ is regular on $Y$ and defines
$\rho_{j,\sigma}$.  Applying $\sigma$ to $q_j(x_j)=t$ proves
\eqref{eq:cover-factorization}.

Each polynomial cover is finite \'{e}tale over $U$.  Its Galois closure, and
then their finite compositum, remains unramified there.  Hence $\pi$ is
finite \'{e}tale over $U$.  Over a simply connected open set $B$, every
connected component of $\pi^{-1}(B)$ maps biholomorphically to $B$.  The
$\deg q_j$ embeddings give distinct holomorphic roots of $q_j(X)-t$ on
that component and therefore exhaust the inverse branches.
\end{proof}

\begin{theorem}[Functional relations are formal]
\label{thm:functional-relations-formal}
Under the hypotheses of \cref{thm:transfer}, if
\[
  \sum_{j=1}^{N}a_jI_{q_j,s_j,\Delta_j}(z)\equiv0,
\]
then
\begin{equation}\label{eq:formal-functional-relation}
  \sum_{j=1}^{N}a_j\mathsf P(q_j,s_j,\Delta_j)=0
  \qquad\text{in }\Pcal_{\Aone}^{\mathrm{KZ,cpt}}.
\end{equation}
\end{theorem}

\begin{proof}
Choose a common tree and graph chains $\Gamma_j$ as in
\cref{lem:common-tree}.  By the converse part of
\cref{thm:combined-moment-rigidity}, every combined density $\Psi_e$
vanishes.

For an oriented edge $e$, let $\gamma_{j,e,i}$ be the lifted branch chain
corresponding to $x_{j,e,i}$.  Then
\[
  \Gamma_j=\sum_{e,i}f_{j,e,i}\gamma_{j,e,i}.
\]
The relative class on $\Aone$ is determined by its boundary, so bilinearity
gives
\begin{equation}\label{eq:edge-decomposition-formal}
  \mathsf P(q_j,s_j,\Delta_j)
  =\sum_{e,i}f_{j,e,i}
    \mathsf P(q_j,s_j,\partial\gamma_{j,e,i}).
\end{equation}

Take the simultaneous cover $\pi:Y\to\Aone_t$ from
\cref{lem:simultaneous-cover}.  Fix an edge $e$.  Choose a thin simply
connected neighborhood $B_e$ of its interior containing no critical value,
and a component $\widetilde B_e$ of $\pi^{-1}(B_e)$.  We justify that
the lifted edge defines an admissible relative chain.  Since $\pi$ is
finite, it is proper, so the lift of the edge interior has compact closure
in $\pi^{-1}(\overline e)$.  At a point over an endpoint $t_0$, a local
holomorphic coordinate $u$ gives
\[
  t-t_0=u^r
\]
for the ramification index $r\ge1$.  This model shows that the lift has a
unique limit at each endpoint.  Reparametrizing a terminal straight
segment by an $r$-th power makes its lift $C^1$ up to that endpoint.
Doing this at both ends gives a finite piecewise $C^1$ chain $\eta_e$.
Its endpoints lie in finite fibers over the algebraic vertices of $T$;
since $\pi$ is defined over $\Qbar$, these endpoints belong to
$Y(\Qbar)$.  Thus $[\eta_e]\in H_1^{\mathrm{rel}}(Y)$.  Orient it so
that $\pi_*\eta_e=\overline e$ with the chosen orientation.

The morphisms from \cref{lem:simultaneous-cover}, restricted to
$\widetilde B_e$, give all inverse branches.  Relabeling them, we obtain
morphisms $\rho_{j,e,i}:Y\to\Aone_x$ such that
\begin{equation}\label{eq:edge-pushforward}
  \rho_{j,e,i*}\eta_e=\gamma_{j,e,i}.
\end{equation}
Define the global regular differential
\begin{equation}\label{eq:Omega-e}
  \Omega_e
  =\sum_{j,i}a_jf_{j,e,i}
    \rho_{j,e,i}^*(s_j(x)\,dx)
  \qquad\text{on }Y.
\end{equation}
On the interior of $\eta_e$, with $t=\pi(y)$, one has
\[
  \Omega_e
  =\left(
      \sum_{j,i}a_jf_{j,e,i}
       \frac{s_j(x_{j,e,i}(t))}{q_j'(x_{j,e,i}(t))}
    \right)dt
  =\Psi_e(t)\,dt=0.
\]
A regular differential on a connected smooth curve that vanishes on a
nonempty analytic arc vanishes identically.  Therefore $\Omega_e=0$ on
$Y$.

Functoriality, \eqref{eq:cover-factorization},
\eqref{eq:edge-pushforward}, and bilinearity now give
\begin{align*}
  \sum_{j,i}a_jf_{j,e,i}
    \mathsf P(q_j,s_j,\partial\gamma_{j,e,i})
  &=\left\langle
      Y,\pi,\Omega_e,[\eta_e]
    \right\rangle\\
  &=0.
\end{align*}
Each displayed identity is a separate relative-chain relation, so the
lifts $\eta_e$ may be chosen independently for different edges.  Summing
these identities and using \eqref{eq:edge-decomposition-formal} proves
\eqref{eq:formal-functional-relation}.
\end{proof}

\section{Proof of the formal-period theorem and consequences}
\label{sec:main-proof}

\begin{proof}[Proof of \cref{thm:main}]
Let an element in the kernel of the period map be represented by
\begin{equation}\label{eq:general-kernel-element}
  \sum_{j=1}^{N}a_j\mathsf P(q_j,r_j,\Delta_j)
  +\sum_{\nu=1}^{M}b_\nu\mathsf E(\beta_\nu).
\end{equation}
Use \cref{lem:small-degree} to replace every term with $\deg q_j\le1$ by
zero-dimensional symbols.  For every remaining term, apply
\cref{lem:twisted-reduction} and collect equal exponents.  The
kernel element is formally equal to
\begin{equation}\label{eq:reduced-kernel-element}
  \sum_{j=1}^{N'}a_j\mathsf P(q_j,s_j,\Delta_j)
  +\sum_{\nu=1}^{M'}c_\nu\mathsf E(\gamma_\nu),
\end{equation}
where $\deg q_j\ge2$, $\deg s_j\le\deg q_j-2$, and the $\gamma_\nu$ are
distinct.

Its numerical value is zero.  If $N'=0$, the
Lindemann--Weierstrass theorem immediately gives
$c_1=\cdots=c_{M'}=0$.  If $N'>0$, \cref{thm:transfer} gives
\[
  \sum_{j=1}^{N'}a_jI_{q_j,s_j,\Delta_j}(z)\equiv0
  \qquad\text{and}\qquad
  c_1=\cdots=c_{M'}=0.
\]
In the latter case the functional relation is zero in the formal period
space by \cref{thm:functional-relations-formal}; in either case the
exponential sum is zero by ordinary linearity.  Hence \eqref{eq:reduced-kernel-element}, and
therefore \eqref{eq:general-kernel-element}, is zero.
\end{proof}

\subsection{Algebraic values with polynomial differentials}

\begin{theorem}[Algebraic-value normal form]
\label{thm:algebraic-value-normal-form}
Let $q\in\Qbar[x]$ have degree $n\ge2$, let
$\Delta=\sum_\alpha d_\alpha[\alpha]$ be a finite degree-zero zero-cycle,
and let $r\in\Qbar[x]$.  Write uniquely
\[
  r=R'+q'R+s,
  \qquad \deg s\le n-2,
\]
and put
\[
  c_\lambda
  =\sum_{\substack{\alpha\\q(\alpha)=\lambda}}
    d_\alpha R(\alpha)
  \qquad(\lambda\in\Qbar).
\]
Then
\[
  I_{q,r,\Delta}(1)\in\Qbar
\]
if and only if
\begin{equation}\label{eq:normal-form-conditions}
  I_{q,s,\Delta}(z)\equiv0
  \qquad\text{and}\qquad
  c_\lambda=0\quad(\lambda\ne0).
\end{equation}
In that case
\begin{equation}\label{eq:normal-form-value}
  I_{q,r,\Delta}(1)=c_0.
\end{equation}
\end{theorem}

\begin{proof}
Twisted reduction gives
\[
  I_{q,r,\Delta}(1)
  =I_{q,s,\Delta}(1)+\sum_\lambda c_\lambda e^\lambda.
\]
If the left side is an algebraic number $a$, regard $-a$ as the coefficient
of $e^0$ and apply \cref{thm:transfer} after collecting equal exponents.
The reduced function must vanish identically, and
Lindemann--Weierstrass gives $c_\lambda=0$ for $\lambda\ne0$ and
$c_0=a$.  The converse follows immediately from the displayed reduction.
\end{proof}

For $r=1$ and $\deg q\ge2$, one has $R=0$ and $s=1$, so
\cref{thm:algebraic-value-normal-form} recovers the algebraic-value rigidity
of Paper~I.  The linear case is already contained in
\cref{lem:small-degree}.  Thus polynomial differentials add precisely the
boundary terms forced by twisted Stokes.

\subsection{A terminal-value independence criterion}

\begin{proposition}[Terminal-value criterion]
\label{prop:terminal-value-criterion}
For $1\le j\le N$, let
$F_j(z)=I_{q_j,s_j,\Delta_j}(z)$ satisfy the reduced hypotheses of
\cref{thm:transfer}.  Suppose that for every $j$ there is an endpoint
$\alpha_j\in\supp\Delta_j$ such that:
\begin{enumerate}[label=\textup{(\roman*)},leftmargin=2.4em]
\item the coefficient $d_{j,\alpha_j}$ of $[\alpha_j]$ in $\Delta_j$ is
nonzero;
\item $q_j'(\alpha_j)s_j(\alpha_j)\ne0$;
\item $\lambda_j=q_j(\alpha_j)$ is not a critical value of any $q_k$, and
is not equal to $q_k(\alpha)$ for any
$\alpha\in\supp\Delta_k$ except the designated endpoint
$(k,\alpha)=(j,\alpha_j)$.
\end{enumerate}
Then $F_1,\ldots,F_N$ are linearly independent over $\Qbar$ as entire
functions.  Consequently their values at $1$ remain linearly independent
after adjoining any finite family of zero-dimensional exponential periods with distinct algebraic
exponents.
\end{proposition}

\begin{proof}
Choose the common tree so that each $\lambda_j$ is the terminal vertex of
its own arm, and construct the graph chains using simple paths.  Let $e_j$
be the terminal open edge of this arm.  Since $\lambda_j$ is a regular value
of every $q_k$, each lift of the closed arm is a leaf edge of
$q_k^{-1}(T)$.

For $k=j$, the leaf ending at $\alpha_j$ occurs in the simple path from the
base vertex to $\alpha_j$ with coefficient $d_{j,\alpha_j}$.  It cannot
occur in a simple path to any other endpoint, because entering and leaving a
leaf would repeat its attachment vertex.  By (iii), no graph chain from any
other datum uses a leaf above the same terminal arm.  Hence the combined
density on $e_j$ for a relation $\sum_k a_kF_k$ is
\[
  a_jd_{j,\alpha_j}
  \frac{s_j(x_j(t))}{q_j'(x_j(t))},
\]
where $x_j(t)$ is the branch ending at $\alpha_j$.

If $\sum_k a_kF_k\equiv0$, then
\cref{thm:combined-moment-rigidity} makes this density identically zero.
Conditions (i) and (ii) imply $a_j=0$.  This holds for every $j$, proving
functional independence.  The numerical assertion follows from
\cref{cor:independence-with-points}.
\end{proof}

\begin{corollary}[An explicit family]
\label{cor:explicit-family}
For every $m\ge1$, the numbers
\[
  1,
  \quad
  \int_0^1e^{x^2}\,dx,
  \quad
  \int_0^1e^{2x^2}\,dx,
  \quad\ldots,\quad
  \int_0^1e^{mx^2}\,dx
\]
are linearly independent over $\Qbar$.  They remain linearly independent
after adjoining finitely many $e^{\beta_1},\ldots,e^{\beta_s}$ with
distinct nonzero algebraic exponents.
\end{corollary}

\begin{proof}
Apply \cref{prop:terminal-value-criterion} with
$q_j(x)=jx^2$, $s_j=1$, and $\Delta_j=[1]-[0]$.  The designated terminal
values are the distinct regular values $j$.  The zero-dimensional exponential periods are handled
by \cref{cor:independence-with-points}.
\end{proof}

\section{Scope of the theorem}
\label{sec:scope}

The word \emph{linear} refers to relations among period values, not to the
degree of the polynomial.  The theorem permits arbitrary polynomial
regular functions and polynomial differentials on the affine line, but its
scope has three essential boundaries.

First, the chains are compact and the source is $\Aone$.  For a finite map
from a general algebraic curve, the Taylor coefficients
\[
  \int_\Gamma f^m\omega
\]
are generally ordinary one-periods rather than algebraic numbers.  Their
linear arithmetic is governed by the theory of one-periods
\cite{HuberWustholz}, but Beukers' $E$-function specialization theorem no
longer applies directly.

Second, products of the periods are higher-dimensional exponential periods:
\[
  \left(\int_{\Gamma_1}e^{q_1}\omega_1\right)
  \left(\int_{\Gamma_2}e^{q_2}\omega_2\right)
  =\int_{\Gamma_1\times\Gamma_2}
     e^{q_1\boxplus q_2}\,\omega_1\boxtimes\omega_2.
\]
Thus algebraic relations of degree greater than one require a
higher-dimensional linear theorem.  On the differential-equation side,
tensor constructions widen the residue spectrum and may introduce integral
differences, so \cref{thm:direct-sum-rigidity} does not extend formally.

Third, noncompact rapid-decay chains require the homology of irregular
connections.  The perfect de Rham--rapid-decay pairing is available in broad
generality \cite{BlochEsnault,HienFlat}, but even the elementary integrals
\[
  \int_0^\infty x^m e^{-x^n}\,dx
  =\frac1n\Gamma\!\left(\frac{m+1}{n}\right)
\]
lead to unresolved algebraic-relation questions for Gamma values.  No
arithmetic specialization theorem for this noncompact setting is used here.

Finally, allowing arbitrary dimensions and setting the regular function in
the exponential equal to zero recovers ordinary periods.  The unrestricted
exponential-period conjecture therefore contains the ordinary
Kontsevich--Zagier period conjecture.  The present theorem concerns exactly
the compact affine-line subspace of \cref{def:affine-line-space}.

\section*{Acknowledgments and AI-assisted research disclosure}

This manuscript is a companion to \cite{LongEIT}.  During its preparation
the author used OpenAI ChatGPT 5.6 Pro and Anthropic Claude Opus 5 for proof
exploration, adversarial checking, symbolic and notational consistency,
bibliographic verification, and language editing.  The AI systems are not
authors.  The authors reviewed the arguments and bear full
responsibility for every statement, proof, citation, and remaining error.

\small
\begin{thebibliography}{99}

\bibitem{AdamczewskiRivoal}
B.~Adamczewski and T.~Rivoal,
\emph{Exceptional values of $E$-functions at algebraic points},
Bull. Lond. Math. Soc. \textbf{50} (2018), no.~4, 697--708.
\href{https://doi.org/10.1112/blms.12168}{doi:10.1112/blms.12168}.

\bibitem{Baker}
A.~Baker,
\emph{Transcendental Number Theory},
Cambridge Mathematical Library,
Cambridge University Press, Cambridge, 1990.

\bibitem{Beukers}
F.~Beukers,
\emph{A refined version of the Siegel--Shidlovskii theorem},
Ann. of Math. (2) \textbf{163} (2006), no.~1, 369--379.
\href{https://doi.org/10.4007/annals.2006.163.369}{doi:10.4007/annals.2006.163.369}.

\bibitem{BlochEsnault}
S.~Bloch and H.~Esnault,
\emph{Homology for irregular connections},
J. Th\'eor. Nombres Bordeaux \textbf{16} (2004), no.~2, 357--371.
\href{https://doi.org/10.5802/jtnb.450}{doi:10.5802/jtnb.450}.

\bibitem{BostanRivoalSalvy}
A.~Bostan, T.~Rivoal, and B.~Salvy,
\emph{Minimization of differential equations and algebraic values of
$E$-functions},
Math. Comp. \textbf{93} (2024), no.~347, 1427--1472.
\href{https://doi.org/10.1090/mcom/3912}{doi:10.1090/mcom/3912}.

\bibitem{CommelinHabeggerHuber}
J.~Commelin, P.~Habegger, and A.~Huber,
\emph{Exponential periods and o-minimality},
arXiv:2007.08280v4 (2025), to appear in
Ann. Inst. Fourier (Grenoble).

\bibitem{CressonViuSos}
J.~Cresson and J.~Viu-Sos,
\emph{On the equality of periods of Kontsevich--Zagier},
J. Th\'eor. Nombres Bordeaux \textbf{34} (2022), no.~2, 323--343.
\href{https://doi.org/10.5802/jtnb.1204}{doi:10.5802/jtnb.1204}.

\bibitem{Delaygue}
\'E.~Delaygue,
\emph{A Lindemann--Weierstrass theorem for $E$-functions},
J. Reine Angew. Math. \textbf{820} (2025), 75--85.
\href{https://doi.org/10.1515/crelle-2024-0090}{doi:10.1515/crelle-2024-0090}.

\bibitem{FischlerRivoalEffective}
S.~Fischler and T.~Rivoal,
\emph{Effective algebraic independence of values of $E$-functions},
Math. Z. \textbf{305} (2023), article~48.
\href{https://doi.org/10.1007/s00209-023-03373-9}{doi:10.1007/s00209-023-03373-9}.

\bibitem{FresanJossen}
J.~Fres\'an and P.~Jossen,
\emph{A non-hypergeometric $E$-function},
Ann. of Math. (2) \textbf{194} (2021), no.~3, 903--942.
\href{https://doi.org/10.4007/annals.2021.194.3.7}{doi:10.4007/annals.2021.194.3.7}.

\bibitem{GavrilovPakovich}
L.~Gavrilov and F.~Pakovich,
\emph{Moments on Riemann surfaces and hyperelliptic Abelian integrals},
Comment. Math. Helv. \textbf{89} (2014), no.~1, 125--155.
\href{https://doi.org/10.4171/CMH/314}{doi:10.4171/CMH/314}.

\bibitem{HienFlat}
M.~Hien,
\emph{Periods for flat algebraic connections},
Invent. Math. \textbf{178} (2009), no.~1, 1--22.
\href{https://doi.org/10.1007/s00222-009-0185-7}{doi:10.1007/\allowbreak s00222-009-0185-7}.

\bibitem{HuberMullerStach}
A.~Huber and S.~M\"uller-Stach,
\emph{Periods and Nori Motives},
Ergebnisse der Mathematik und ihrer Grenzgebiete, 3.~Folge, vol.~65,
Springer, Cham, 2017.
\href{https://doi.org/10.1007/978-3-319-50926-6}{doi:10.1007/978-3-319-50926-6}.

\bibitem{HuberWustholz}
A.~Huber and G.~W\"ustholz,
\emph{Transcendence and Linear Relations of 1-Periods},
Cambridge Tracts in Mathematics, vol.~227,
Cambridge University Press, Cambridge, 2022.
\href{https://doi.org/10.1017/9781009019729}{doi:10.1017/9781009019729}.

\bibitem{KontsevichZagier}
M.~Kontsevich and D.~Zagier,
\emph{Periods},
in \emph{Mathematics Unlimited--2001 and Beyond},
Springer, Berlin, 2001, pp.~771--808.
\href{https://doi.org/10.1007/978-3-642-56478-9_39}{doi:10.1007/978-3-642-56478-9\_39}.

\bibitem{LongEIT}
C.~D.~Long and A.-A.~Le~Blanc,
\emph{Algebraic Exponential Integrals: Transcendence and Linear Relations},
companion preprint, August 2026.

\bibitem{PakovichMuzychuk}
F.~Pakovich and M.~Muzychuk,
\emph{Solution of the polynomial moment problem},
Proc. Lond. Math. Soc. (3) \textbf{99} (2009), no.~3, 633--657.
\href{https://doi.org/10.1112/plms/pdp010}{doi:10.1112/plms/pdp010}.

\bibitem{PakovichRoytvarfYomdin}
F.~Pakovich, N.~Roytvarf, and Y.~Yomdin,
\emph{Cauchy-type integrals of algebraic functions},
Israel J. Math. \textbf{144} (2004), 221--291.
\href{https://doi.org/10.1007/BF02916714}{doi:10.1007/BF02916714}.

\end{thebibliography}
\end{document}
